\chapter*{Résumé}
\phantomsection
\addcontentsline{toc}{chapter}{Résumé}

Ce rapport présente un projet de fin d'études réalisé au sein de STMicroelectronics Tunis
afin d'étudier l'écosystème STM32Cube face à des écosystèmes concurrents. Dans une mission
indépendante confiée au stagiaire, nous avons d'abord recherché les cartes STM32 les plus
proches d'une liste de cartes concurrentes fournie par l'équipe d'accueil. Cette étude de
positionnement produit n'a pas déterminé les plateformes des benchmarks ultérieurs. Le
travail expérimental a ensuite commencé par des pilotes USB HID, CDC et MSC comparant STM32
à Renesas et NXP. L'optimisation de la référence STM32 a réduit la RAM et le temps
d'initialisation dans les trois scénarios Renesas. Face à NXP, l'interprétation de la ROM
dépend de l'inclusion du RCC, tandis que MCX-N9 conserve les valeurs rapportées de RAM et
d'initialisation les plus faibles.

Nous avons établi une méthode de benchmarking fondée sur des
fonctionnalités équivalentes, des conditions de compilation documentées, l'empreinte mémoire
extraite des fichiers linker map, des mesures d'exécution lorsqu'elles étaient étayées par
des observations matérielles, ainsi qu'une vérification explicite de l'équivalence sémantique
et de l'équité des comparaisons. Nous avons également conçu deux outils de support~:
\emph{MCU Workflow Studio}, qui structure l'analyse inter-vendeurs des fichiers linker map,
et le \emph{MCU Benchmark Copilot Agent}, qui formalise le processus de création, de portage,
de compilation, de traçabilité et de validation des projets.

La campagne ST--GD32 achevée couvre la pile USB device à travers les scénarios Human
Interface Device (HID) et Communication Device Class (CDC), ainsi que FreeRTOS à travers des
scénarios d'ordonnancement basique, coopératif et préemptif. STM32 nécessite 35--88\,\% de
mémoire flash supplémentaire selon les scénarios, principalement en raison de la couche
Hardware Abstraction Layer (HAL), mais utilise 57--65\,\% de RAM en moins dans les scénarios
USB. Dans les observations rapportées, le temps d'initialisation USB est inférieur de
61,7\,\% sur
STM32, les mesures USB en régime établi sont proches et les mesures FreeRTOS diffèrent au
maximum de 3,3\,\%. Ces résultats nous ont conduits à retenir le module Reset and Clock
Control (RCC) de la HAL et la sélection des fonctionnalités USB à la compilation comme pistes
d'investigation contrôlée, et non comme optimisations déjà démontrées.

La comparaison ST--Texas Instruments associe une analyse statique des pilotes Low-Layer de
STM32CubeC5 et de la bibliothèque de pilotes MSPM33 à douze projets bare-metal appariés
couvrant la conversion analogique-numérique, les entrées-sorties généralistes,
l'I\textsuperscript{2}C, le SPI et l'UART. Sur l'ensemble de ces douze images liées
indépendamment, STM32C5 utilise 2,8\,\% de mémoire flash agrégée supplémentaire et 5,0\,\% de
RAM agrégée en moins. Ces résultats caractérisent la structure du code, la couverture des
Software Development Kits (SDK) et l'empreinte des images complètes. Nous avons ensuite
comparé FreeRTOS à travers des scénarios basique, coopératif et préemptif. Les trois images
STM32C5 utilisent 27,9--29,6\,\% de flash en moins, tandis que MSPM33 économise 56 à 76
octets de RAM. Des observations d'exécution ont également été collectées, mais STM32C5
fonctionne à 144\,MHz contre 32\,MHz pour MSPM33. Après normalisation par la fréquence, le
scénario basique devient presque équivalent; les scénarios coopératif et préemptif restent
affectés par la configuration du noyau et par des primitives de synchronisation différentes.
Ces limites empêchent de conclure globalement sur l'efficacité de l'ordonnanceur.

\noindent\rule[2pt]{\textwidth}{0.5pt}

{\textbf{Mots clés :}}
STM32Cube, benchmarking, empreinte firmware, microcontrôleur, FreeRTOS, USB, GD32, Texas Instruments, pilotes bas niveau, comparaison inter-vendeurs.
\\
\noindent\rule[2pt]{\textwidth}{0.5pt}